%%%%%%%%%%Fourth Chapter%%%%%%%%%%
\chapter{Simulation Results and Discussion}

\section{Overview}
This chapter presents the simulation results obtained from the metasurface transmitarray lens antenna design and the communication channel evaluation. The results are organized into antenna-level electromagnetic performance and communication-level channel performance. The purpose of this chapter is to evaluate whether the designed antenna can support the required communication link.

\section{Unit Cell Simulation Results}
The unit cell simulation results are used to verify the transmission magnitude and phase response of the selected metasurface element. The unit cell should provide sufficient phase coverage while maintaining low transmission loss across the target operating frequency.

\subsection{Transmission Magnitude}
The transmission magnitude result shows how much electromagnetic power passes through the unit cell. A high transmission magnitude is desired because it indicates lower insertion loss and better aperture efficiency.

\subsection{Transmission Phase}
The transmission phase result is used to map the required lens phase compensation to the corresponding unit cell geometry. The phase response should cover the required phase range with smooth variation between design states.

\begin{table}[h!]
	\begin{center}
		\caption{Summary of unit cell simulation results}
		\label{table:unit_cell_results}
		\begin{tabular}{|l|c|c|}
			\hline
			\textbf{Parameter} & \textbf{Target/Condition} & \textbf{Simulation Result} \\
			\hline
			Operating frequency & TBD & TBD \\
			Transmission magnitude & TBD & TBD \\
			Transmission phase range & TBD & TBD \\
			Unit cell size & TBD & TBD \\
			Bandwidth & TBD & TBD \\
			\hline
		\end{tabular}
	\end{center}
\end{table}

\section{Transmitarray Lens Simulation Results}
The full transmitarray lens simulation evaluates the antenna radiation performance after the unit cells are arranged over the aperture. This analysis shows whether the lens can transform the feed radiation into the desired directional beam.

\subsection{Reflection Coefficient}
The reflection coefficient is evaluated to observe impedance matching near the operating frequency. A lower reflection coefficient indicates that more input power is accepted by the antenna system.

\subsection{Radiation Pattern}
The radiation pattern shows the angular distribution of radiated power. The main beam direction, beamwidth, and side-lobe level are observed to evaluate the focusing and beam-forming performance of the transmitarray lens.

\subsection{Gain and Directivity}
Gain and directivity are used to evaluate the ability of the transmitarray lens antenna to concentrate electromagnetic energy in the desired direction. A higher gain indicates better suitability for directional communication links.

\subsection{Aperture Efficiency}
Aperture efficiency describes how effectively the physical aperture is converted into radiated gain. This parameter is influenced by transmission loss, phase error, feed illumination, spillover, and aperture taper.

\begin{table}[h!]
	\begin{center}
		\caption{Summary of transmitarray antenna simulation results}
		\label{table:antenna_results}
		\begin{tabular}{|l|c|}
			\hline
			\textbf{Parameter} & \textbf{Simulation Result} \\
			\hline
			Reflection coefficient & TBD \\
			Peak gain & TBD \\
			Directivity & TBD \\
			Main beam direction & TBD \\
			Half-power beamwidth & TBD \\
			Side-lobe level & TBD \\
			Aperture efficiency & TBD \\
			\hline
		\end{tabular}
	\end{center}
\end{table}

\section{Communication Channel Evaluation Results}
The communication channel evaluation connects the antenna simulation results with link-level performance. The antenna radiation pattern and gain are used to evaluate how the designed antenna affects received power, path loss, coverage, and link quality.

\subsection{Evaluation Scenario}
The evaluation scenario includes the transmitter location, receiver location, antenna orientation, propagation distance, and environmental assumptions. These conditions are used to interpret the channel results consistently.

\subsection{Received Power}
Received power is used to evaluate how much signal power is available at the receiver. The result is affected by transmit power, antenna gain, propagation loss, antenna alignment, and multipath propagation.

\subsection{Path Loss}
Path loss describes the reduction of signal power between transmitter and receiver. The evaluated path loss is compared across the selected propagation condition to determine whether the antenna provides sufficient link margin.

\subsection{Signal-to-Noise Ratio}
Signal-to-noise ratio is used to estimate the reliability of the communication link. A higher signal-to-noise ratio indicates better received signal quality and improved communication performance.

\subsection{Coverage Evaluation}
Coverage evaluation is used to observe the spatial region where the communication link can satisfy the required performance. The coverage result can be used to compare antenna configurations, beam directions, or propagation scenarios.

\begin{table}[h!]
	\begin{center}
		\caption{Summary of communication channel evaluation results}
		\label{table:channel_results}
		\begin{tabular}{|l|c|}
			\hline
			\textbf{Parameter} & \textbf{Evaluation Result} \\
			\hline
			Propagation distance & TBD \\
			Received power & TBD \\
			Path loss & TBD \\
			Signal-to-noise ratio & TBD \\
			Delay spread & TBD \\
			Coverage area & TBD \\
			\hline
		\end{tabular}
	\end{center}
\end{table}

\section{Comparison and Discussion}
The simulation results are discussed by comparing antenna performance and channel performance. The relationship between antenna gain, radiation direction, path loss, and received power is analyzed to understand the impact of the metasurface transmitarray lens antenna on the communication channel.

\section{Summary}
This chapter summarizes the main simulation findings. The unit cell results verify the phase and transmission behavior of the metasurface element. The antenna results evaluate radiation performance, while the communication channel results show the link-level effect of the designed antenna.
